% Generated from common/cover-notes.json by tools/series_cover.py.
\backmatter
\cleardoublepage
\chapter*{About the Cover}
\addcontentsline{toc}{chapter}{About the Cover}
\noindent Every reduced fraction carries a circle whose radius is one over twice the square of its denominator, standing on the axis and tangent to its neighbours. The pale arcs join fractions that are adjacent in the Farey sequence. Lester Ford published the circles in the American Mathematical Monthly in 1938, in a note called Fractions. Their tangencies are the geometry behind the Farey sequence studied in the chapter Arithmetic Functions on Average, where the spacing of those fractions is tied to the Riemann hypothesis. The drawing was computed by the book's own TikZ cover generator.
